\documentclass[10pt]{article}

\usepackage[utf8]{inputenc}

\usepackage[T1]{fontenc}

\usepackage{amsmath}

\usepackage{amsfonts}

\usepackage{amssymb}

\usepackage[version=4]{mhchem}

\usepackage{stmaryrd}

\usepackage{graphicx}

\usepackage[export]{adjustbox}

\graphicspath{ {./images/} }

\usepackage{bbold}



\begin{document}

ного представителя $x$, к нахождению стабилизатора $G_{x}$ и к описанию какого-либо - по возможности обозримого - класса функций, постоянных на О. (инвариантов) и разделяющих разные О.; эти функции позволяют описать «расположение» 0 . в $X$ (О. являются пересечениями их множеств уровня). Эта программа обычно наз.



\includegraphics[max width=\textwidth, center]{2023_07_08_bc0bf64cb896270e50b3g-1}

К такой задаче часто сводятся многие задачи классификации. Так, в примере 2) это - задача классификации билинейных симметрич. форм с точностью до эквивалентности; инварианты в әтом случае «дискретны»это ранг и сигнатура, а стабилизатор $G_{f}$, где $f$ невырождена, является соответствующей псевдоортогональной группой. Классич. теория жордановой формы матриц (также, как и теории других нормальных борм матриц) тоже укладывается в эту схему: жорданова форма - это канонич. представитель (определенный, правда, с точностью до порядка жордановых клеток) в $\mathrm{O}$. полной линейной группы $\mathrm{GL}_{n}(\mathbb{C})$ на пространстве всех комплексных $(n \times n)$-матриц при действии, заданном формулой $Y \rightarrow A Y A^{-1}$; среди инвариантов важное место занимают коәфффициенты характеристич. многочлена матрицы $Y$ (к-рые не разделяют, однако, любые две О.). Идея рассмотрения әквивалентных объектов как О. нек-рой группы активно используется в различных задачах классификации, напр. в алгебрапч. жодулей теории (теория Мамфорда, см. [10]), в теории перечисления графов (см. [2]) и др.



Если $G$ и $X$ конечны, то



\[

|X / G|=\frac{1}{|G|} \sum_{g \in G} \mid \text { Fix } g \mid

\]



где $|Y|$ - число элементов множества $Y$, а



\[

\text { Fix } g=\{x \in X \mid g(x)=x\} \text {. }

\]



Если $G$ - компактная группа Ли, гладко действующая на связном гладком многообразии $X$, то структура О. на $X$ локально конечна, т. е. у любой точки $x \in X$ существует такая окрестность $U$, что число различных стабилизаторов $G_{y}, y \in U$, с точностью до сопряженности в $G$ конечно. В частности, если $X$ компактно, то конечно число различных (с точностью до сопряженности в $G$ ) стабилизаторов $G_{y}, y \in X$. При этом для любой подгруппы $H$ в $G$ каждое из множеств



\[

X_{(H)}=\left\{x \in X \cdot \mid G_{x} \text { сопряжена } H \text { в } G\right\}

\]



является пересечением открытого и замкнутого инвариантных подмножеств в $X$. Исследование $X_{(H)}$ приводит в этом случае к классификации действий (см. [1]). Аналоги этих результатов получены в геометрич. инвариантов теории (см. [3]). А именно, пусть $G$ - редуктивная алгебраич. группа, регулярно действующая на аффинном алгебраич. многообразии $X$ (основное поле $k$ алгебраически замкнуто и имеет характеристику 0). В замыкании любой О. содержится единствөнная замкнутая 0 . Существует разбиение $X$ в объединение конечного числа локально замкнутых инвариантных непересекающихся подмножеств $X=\mathrm{U}_{\alpha} X_{\alpha}$, обладающее своїствами: а) если $x, y \in X_{\alpha}$ и $G(x)$ замкнута, то стабилизатор $G_{y}$ сопряжен в $G$ подгруппе в $G_{x}$, а если замкнута и $G(y)$, то $G_{y}$ сопряжен с $G_{x} ;$ б) если $x \in X_{\alpha}, y \in X_{\beta}, \alpha \neq \beta$, а $G(x)$ и $G(y)$ замкнуты, то $G_{x}$ и $G_{y}$ не сопряжены в $G$. Если $X$ - гладкое алгебраич. многообразие (напр., в важном случае, когда рассматривается рациональное линейное представление $G$ в векторном пространстве $V=X$, то существует такое непустое открытое подмножество $\Omega$ в $X$, что $G_{x}$ и $G_{y}$ сопряжены в $G$ для любых $x, y \in \Omega$. Последний результат является утверждением о свойстве то ч е к 0 б $е$ го. Поло же ни я в $X$, т. е. точек, заполняющих непустое открытое подмножество; имеется и ряд других утверждений такого типа. Напр., для рационального линейного представления полупростой группы $G$ в векторном пространстве $V 0$. точек общего положения замкнуты тогда и только тогда, когда их стабилизаторы редуктивны (см. [7]); в случае, когда $G$ неприводима, найден явный вид стабилизаторов точек общего положения (см. [8], [9]). Вопрос о замкнутости О. является специфическим и важным в этой теории. Так, множество тех точек $x \in V, 0$. к-рых содержит в своем замыкании вуль пространства $V$, совпадает с многообразием нулей непостоянных инвариантных многочлөнов на $V$; во многих случаях, и в частности в применениях теории инвариантов к теории модулей, это многообразие играет существенную роль (см. [10]). Любые две различные замкнутые 0 . разделяются инвариантными многочленами. Орбита $G(x)$ замкнута тогда и только тогда, когда замкнута О. точки $x$ относительно нормаливатора $G(x)$ в $G$ (см. [4]). Появление незамкнутых О. связано со свойствами $G$; если $G$ унипотентна (а $X$ афффинно), то любая О. замкнута (см. [6]). Одним из направлений теории инвариантов является изучение орбитальных разложений различных конкретных действий (в особенности линейных представлений). Одно из них - присоединенное представление редуктивной группы $G$ - подробно исследовалось (см., напр., [11]). Его изучение связано с теорией представлений группы $G$; см. Орбит метод.



Jum.: [1] P a l a i s R., The classification of G-spaces, Providence, 1960 (Mem. Amer. Math. Soc., N No $^{36)}$ ); [2] X a p a p u $\Phi$. , Math. France. Mém. 33», 1973 , p. 81-105; [4] e r o « e, «Invent. Math.», 1975, v. 29, № 3, p. 231 - 38; [5] Б о р е л ь А., Јинейные алгебраические группы, пер. с англ., M., 1972; [6] s t e i n-



\includegraphics[max width=\textwidth, center]{2023_07_08_bc0bf64cb896270e50b3g-1(3)}

1970, т. 34, с. 523-31; [8] П1 о п о в А. М., "Функц. анализ и его



\includegraphics[max width=\textwidth, center]{2023_07_08_bc0bf64cb896270e50b3g-1(2)}

Tam

metric invariant theory, B. metric invariant theory, B.- Hdlb. - N. Y., 1965; [11] K o s-

t a n t B., "Amer. J. Math.", 1963, v. 85, № 3, p. 327-404; [12] Х а м ф р и Дж., Јинейные алгебраические группы, пер. с англ.,



ОРБИТАЛЬНАЯ УСТОЙЧИВОСТЬ - свойство траектории $\xi$ (решения $x(t)$ ) автономной системы обыкшовенных дифференциальных уравнений



\[

\dot{x}=f(x), x \in \mathbb{R}^{n},

\]



состоящее в следующем: для всякого $\varepsilon>0$ существует $\delta>0$ такое, что всякая положительная полутраектория,



\begin{center}

\includegraphics[max width=\textwidth]{2023_07_08_bc0bf64cb896270e50b3g-1(4)}

\end{center}



\includegraphics[max width=\textwidth, center]{2023_07_08_bc0bf64cb896270e50b3g-1(1)}

$\mathrm{p}$ и е й понимается множество значений решения $x(t), t \in \mathbb{R}$, системы $(*)$, а под п о ло ж и т е л ь н о й п о л у т р а е к т о р и е й - множество значений решения $x(t)$ при $t \geqslant 0$. Если решение $x(t)$ устойчиво по Ляпунову, то его траектория орбитально устойчива.



Траектория $\xi$ наз. а с и м п т о т и ч е с к и о р б ит а л ь о у с т о й ч и в о й, если она орбитально устойчива б, кроме того, найдется $\delta_{0}>0$ такое, что траектория всякого решения $x(t)$ системы $(*)$, начинающегося в $\delta_{0}$-окрестности траектории $\xi$ (т. е. $d(x(0), \xi)<$ $\left.<\delta_{0}\right)$, стремится при $t \rightarrow+\infty$ к траектории $\xi$, то өсть



\[

\lim _{t \rightarrow+\infty} d(x(t), \xi)=0

\]



где



\[

d(x, \xi)=\inf _{y \in \xi} d(x, y)

\]



\begin{itemize}

  \item расстояние от точки $x$ до множества $\xi(d(x, y)-$ расстояние между точками $x$ и $y$ ).

\end{itemize}



Роль понятия асимптотической орбитальной устойчивости основана на следующих фактах. Периодич. решение системы (*) никогда не бывает асимптотически устойчивым. Но если у периодич. решения такой системы модули всех мультипликаторов, кроме одного, меньше единицы, то траектория этого периодич. решения асимптотически орбитально устойчива (А н д р о-





\end{document}